\subsection{Partitions of unity, contractions and absolute retracts}\label{subsec:preliminaries-topology}
\paraid{p-33a116cd3f4a}
We recall partitions of unity,
an extension theorem for continuous real functions,
and the terminology for contractions and absolute retracts.

\paraid{p-4f4ac6ddbc21}
A family of subsets of a topological space is \emph{locally finite} if every point
has a neighborhood meeting only finitely many members of the family.
A cover \emph{refines} another cover if each member of the first is contained
in a member of the second.
A Hausdorff space is \emph{paracompact} if every open cover admits a locally finite
open refinement.

\paraid{p-bcbfe50f69ab}
Let
$X$
be a topological space and let
$I$
be an index set.
A family of continuous functions
$\lambda_i\colon X\to[0,1]$,
indexed by
$i\in I$,
is a \emph{locally finite partition of unity} if the supports
\[
 \supp\lambda_i=\overline{\{x\in X\mid\lambda_i(x)\ne0\}}
\]
form a locally finite family and for every
$x\in X$
we have
\[
 \sum_{i\in I}\lambda_i(x)=1.
\]
It is \emph{subordinate} to an open cover if each support is contained in a member
of that cover.

\begin{theorem}[{\cite[Theorems~1.3.33 and~1.3.38]{CobzasMiculescuNicolae2019}}]\label{thm:metric-partition-unity}
\paraid{p-e8f70dfc07eb}
Every metrizable space is paracompact.
Every open cover of a metrizable space admits a locally finite partition of unity
subordinate to it.
\end{theorem}

\begin{theorem}[{Tietze--Urysohn, \cite[Theorem~2.1.8]{Engelking1989}}]\label{thm:tietze}
\paraid{p-fb8cfff8e576}
Let
$\ExtensionDomain$
be a normal space,
let
$\ClosedDomain\subset\ExtensionDomain$
be closed,
and let
$\FunctionBound\geq0$.
Then every continuous function
$\Mapping\colon\ClosedDomain\to[-\FunctionBound,\FunctionBound]$
has a continuous extension
$\TietzeExtension\colon\ExtensionDomain\to[-\FunctionBound,\FunctionBound]$.
\end{theorem}

\paraid{p-c07e9f5bb5bf}
A continuous map
$\Mapping \colon \ClosedDomain \to \ExtensionTarget $
is \emph{null-homotopic} if there exist a point
$\ConstantPoint \in \ExtensionTarget $
and a continuous map
$\ContractionHomotopy \colon \ClosedDomain \times\UnitInterval\to \ExtensionTarget $
such that for every
$\FirstScale\in\ClosedDomain$
we have
\[
 \ContractionHomotopy (\FirstScale ,0)=\Mapping (\FirstScale ),
\]
and
\[
 \ContractionHomotopy (\FirstScale ,1)=\ConstantPoint.
\]
A nonempty space is \emph{contractible} if its identity map is null-homotopic.

\paraid{p-457cc7aecb3b}
The following proposition shows that scaling distances defines a contraction.
Write
$\SingletonSpace=\MetricSpace{\{0\}}{0}$
for the one-point metric space.

\begin{proposition}\label{lem:scaling-contraction}
\paraid{p-3f7987feb185}
For compact metric spaces
$\MetricSpace{\BaseCarrier }{\MetricSymbol }$
and
$\MetricSpace{\ComparisonCarrier }{\ComparisonMetric }$
and real numbers
$\FirstScale ,\SecondScale \geq0$,
we have
\begin{equation}\label{eq:scaling}
 \begin{aligned}
 \GHDistance(\MetricQuotient{\BaseCarrier }{\FirstScale  \MetricSymbol },\MetricQuotient{\ComparisonCarrier }{\SecondScale  \ComparisonMetric })
 &\leq \FirstScale \GHDistance(\MetricSpace{\BaseCarrier }{\MetricSymbol },\MetricSpace{\ComparisonCarrier }{\ComparisonMetric })
       +\frac{|\FirstScale -\SecondScale |}{2}\diam\MetricSpace{\ComparisonCarrier }{\ComparisonMetric }.
 \end{aligned}
\end{equation}
Consequently the map
\[
 \ContractionHomotopy \colon\GHSpace\times\UnitInterval\to\GHSpace
\]
defined by
\[
 \ContractionHomotopy (\MetricSpace{\BaseCarrier }{\MetricSymbol },\HomotopyTime )=\MetricQuotient{\BaseCarrier }{(1-\HomotopyTime )\MetricSymbol }
\]
is continuous and satisfies
\[
 \ContractionHomotopy (\MetricSpace{\BaseCarrier }{\MetricSymbol },0)=\MetricSpace{\BaseCarrier }{\MetricSymbol },
\]
and
\[
 \ContractionHomotopy (\MetricSpace{\BaseCarrier }{\MetricSymbol },1)=\SingletonSpace.
\]
The contraction fixes the one-point space at every time,
so
\[
 \ContractionHomotopy (\SingletonSpace,\HomotopyTime )=\SingletonSpace.
\]
In particular,
$\GHSpace$
is contractible.
\end{proposition}

\begin{proof}
\paraid{p-98b9e8081dac}
By the triangle inequality and the scaling identities in
\cite[arXiv v1, Examples~6.31 and~6.32]{Tuzhilin2020},
we obtain \eqref{eq:scaling},
also when one of the scaling factors is zero.
Since the diameter is continuous on
$\GHSpace$,
the estimate proves continuity of the map
\[
 \GHSpace\times[0,\infty)\to\GHSpace,
 \qquad
 (\MetricSpace{\BaseCarrier}{\MetricSymbol},\FirstScale)
       \longmapsto\MetricQuotient{\BaseCarrier}{\FirstScale\MetricSymbol},
\]
including at
$\FirstScale=0$.
The three identities for
$\ContractionHomotopy $
follow directly from its definition.
\end{proof}

\paraid{p-3e5d9433b03f}
For balls in Gromov--Hausdorff space we use
a different  symbol,
\[
 \GHball(\MetricSpace{\BaseCarrier}{\MetricSymbol},\Radius)
 =\{\MetricSpace{\ComparisonCarrier}{\ComparisonMetric}\in\GHSpace\mid
 \GHDistance(\MetricSpace{\BaseCarrier}{\MetricSymbol},\MetricSpace{\ComparisonCarrier}{\ComparisonMetric})<\Radius\}.
\]
The correspondence formula also implies
$\GHDistance(\MetricSpace{\BaseCarrier }{\MetricSymbol },\SingletonSpace)=\diam\MetricSpace{\BaseCarrier }{\MetricSymbol }/2$.

\paraid{p-263af0d7c090}
We state extension and retraction properties relative to a class of
topological spaces.

\begin{definition}\label{def:extension}
\paraid{p-3d55c72c99aa}
Let
$\TopologicalSpaceClass$
be a class of topological spaces.
A topological space
$\ExtensionTarget$
is an \emph{absolute extensor for $\TopologicalSpaceClass$},
or AE for
$\TopologicalSpaceClass$,
if for every space
$\BaseCarrier\in\TopologicalSpaceClass$,
every closed subset
$\ClosedDomain\subset\BaseCarrier$,
and every continuous map
$\Mapping\colon\ClosedDomain\to\ExtensionTarget$,
there exists a continuous map
$\ContinuousExtension\colon\BaseCarrier\to\ExtensionTarget$
such that
$\ContinuousExtension|_\ClosedDomain=\Mapping$.
It is an \emph{absolute neighborhood extensor for $\TopologicalSpaceClass$},
or ANE for
$\TopologicalSpaceClass$,
if for every space
$\BaseCarrier\in\TopologicalSpaceClass$,
every closed subset
$\ClosedDomain\subset\BaseCarrier$,
and every continuous map
$\Mapping\colon\ClosedDomain\to\ExtensionTarget$,
there exist an open neighborhood
$\OpenNeighborhood$
of
$\ClosedDomain$
in
$\BaseCarrier$
and a continuous map
$\ContinuousExtension\colon\OpenNeighborhood\to\ExtensionTarget$
such that
$\ContinuousExtension|_\ClosedDomain=\Mapping$.
A space
$\ExtensionTarget\in\TopologicalSpaceClass$
is an \emph{absolute retract for $\TopologicalSpaceClass$},
or AR for
$\TopologicalSpaceClass$,
if for every space
$\BaseCarrier\in\TopologicalSpaceClass$
and every embedding
$\FirstEmbedding\colon\ExtensionTarget\to\BaseCarrier$
with closed image,
there exists a continuous map
$\RetractionMap\colon\BaseCarrier\to\ExtensionTarget$
such that
$\RetractionMap\circ\FirstEmbedding=\IdentityMap_\ExtensionTarget$.
A space
$\ExtensionTarget\in\TopologicalSpaceClass$
is an \emph{absolute neighborhood retract for $\TopologicalSpaceClass$},
or ANR for
$\TopologicalSpaceClass$,
if for every space
$\BaseCarrier\in\TopologicalSpaceClass$
and every embedding
$\FirstEmbedding\colon\ExtensionTarget\to\BaseCarrier$
with closed image,
there exist an open neighborhood
$\OpenNeighborhood$
of
$\FirstEmbedding(\ExtensionTarget)$
in
$\BaseCarrier$
and a continuous map
$\RetractionMap\colon\OpenNeighborhood\to\ExtensionTarget$
such that
$\RetractionMap\circ\FirstEmbedding=\IdentityMap_\ExtensionTarget$.
\end{definition}

\paraid{p-cb58f0b73b6e}
Unless otherwise stated,
ANR means ANR for all metrizable spaces,
and the same convention applies to AE,
ANE,
and AR.
Spaces described by these unqualified terms are assumed to be metrizable.

\begin{theorem}[{\cite[Theorem~12.1]{Dugundji1958}}]\label{thm:retract-extensor}
\paraid{p-a393bc2f3507}
For a metrizable target in the category of metrizable spaces,
the AE and AR properties are equivalent,
and the ANE and ANR properties are equivalent.
\end{theorem}

\begin{definition}\label{def:local-contractibility}
\paraid{p-e8a842a80650}
A space
$\ExtensionTarget $
is \emph{locally contractible at}
$\BasePoint \in \ExtensionTarget $
if for every open neighborhood
$\LargerNeighborhood $
of
$\BasePoint $
there is an open neighborhood
$\SmallerNeighborhood \subset \LargerNeighborhood $
of
$\BasePoint $
such that the inclusion
$\SmallerNeighborhood \to \LargerNeighborhood $
is null-homotopic.
\end{definition}

\paraid{p-6fea115b8d45}
The homotopy may be chosen to end at the constant map with value
$\BasePoint $.
Indeed,
evaluating the homotopy at
$\BasePoint $
defines a path from
$\BasePoint $
to the constant value,
and its reversed path can be appended to the homotopy path of every point of
$\SmallerNeighborhood$.
For
$\GHSpace$,
this definition is equivalent to the following condition.
For every
$\MetricSpace{\BaseCarrier }{\MetricSymbol }\in\GHSpace$
and
$\ErrorTolerance >0$,
there exist
$0<\ContractionRadius <\ErrorTolerance $
and a continuous map
\[
 \ContractionHomotopy \colon \GHball(\MetricSpace{\BaseCarrier }{\MetricSymbol },\ContractionRadius )\times\UnitInterval
          \to \GHball(\MetricSpace{\BaseCarrier }{\MetricSymbol },\ErrorTolerance )
\]
such that
$\ContractionHomotopy (\MetricSpace{\ComparisonCarrier }{\ComparisonMetric },0)=\MetricSpace{\ComparisonCarrier }{\ComparisonMetric }$
and
$\ContractionHomotopy (\MetricSpace{\ComparisonCarrier }{\ComparisonMetric },1)=\MetricSpace{\BaseCarrier }{\MetricSymbol }$.
